\documentclass[11pt,a4paper]{article}

\usepackage[T1]{fontenc}
\usepackage[utf8]{inputenc}
\usepackage[margin=2.3cm]{geometry}
\usepackage{lmodern}
\usepackage{hyperref}
\usepackage{booktabs}
\usepackage{array}
\usepackage{enumitem}

\setlist{noitemsep,topsep=4pt}
\setlength{\parindent}{0pt}
\setlength{\parskip}{6pt}

\title{Chess Club Tournament \& Rating Manager\\User Manual}
\author{Yahya El Idrissi\\THGA Bochum -- Introduction to Database Management Systems}
\date{September 2026}

\begin{document}

\maketitle
\newpage

\section{Overview}

The \textbf{Chess Club Tournament \& Rating Manager} is a desktop application for managing players, tournaments, game results and Elo ratings in a small chess club.

The application supports the following tasks:
\begin{itemize}
    \item create players;
    \item create tournaments and their rounds;
    \item register players for tournaments;
    \item record game results;
    \item update Elo ratings automatically;
    \item view rating history;
    \item view tournament standings with Buchholz tiebreaks.
\end{itemize}

The main window contains six tabs: \textbf{Players}, \textbf{Standings}, \textbf{Rating History}, \textbf{Record Game}, \textbf{Register Player}, and \textbf{Create Tournament}.

The usual workflow is:

\begin{center}
\textbf{Create players $\rightarrow$ Create tournament $\rightarrow$ Register players $\rightarrow$ Record games $\rightarrow$ Check standings and rating history}
\end{center}

\section{Setup and Startup}

\subsection{Requirements}

Install the following before running the application from source:
\begin{itemize}
    \item Docker Desktop or Docker Engine with Docker Compose;
    \item Python 3;
    \item Git, if the repository must be cloned.
\end{itemize}

Create a file named \texttt{.env} in the repository root using \texttt{.env.example} as the template. The frontend API key must match the backend \texttt{API\_KEY}.

For a first source-based setup, create the API virtual environment and install the dependencies.

\textbf{Linux / macOS:}
\begin{verbatim}
python3 -m venv api/.venv
api/.venv/bin/pip install -r api/requirements.txt
\end{verbatim}

\textbf{Windows PowerShell:}
\begin{verbatim}
py -m venv api/.venv
.\api\.venv\Scripts\pip install -r api\requirements.txt
\end{verbatim}

\subsection{Linux (Debian / Ubuntu)}

Linux is the main target environment. From the repository root, the easiest startup method is:

\begin{verbatim}
./scripts/run-local.sh
\end{verbatim}

The script starts PostgreSQL, FastAPI and the tkinter frontend.

If the Debian frontend package is installed, the frontend can also be launched with:

\begin{verbatim}
chess-club-manager
\end{verbatim}

In this case, PostgreSQL and the API must already be running.

\subsection{macOS}

On macOS, start the components manually.

\textbf{Terminal 1:}
\begin{verbatim}
cd /path/to/chess-club-manager
docker compose up -d postgres
set -a
source .env
set +a
api/.venv/bin/uvicorn app.main:app \
  --app-dir api --host 127.0.0.1 --port 8000
\end{verbatim}

Keep this terminal open. Then open a second terminal:

\begin{verbatim}
cd /path/to/chess-club-manager
export CHESS_API_URL="http://127.0.0.1:8000"
export CHESS_API_KEY="$(grep '^API_KEY=' .env | cut -d= -f2-)"
PYTHONPATH="$PWD/frontend" python3 -m chess_club_fe.main
\end{verbatim}

\subsection{Windows}

On Windows, use PowerShell and Docker Desktop.

\textbf{PowerShell 1:}
\begin{verbatim}
cd C:\path\to\chess-club-manager
docker compose up -d postgres
$env:POSTGRES_USER="chess"
$env:POSTGRES_PASSWORD="change-me"
$env:POSTGRES_DB="chessdb"
$env:POSTGRES_PORT="5432"
$env:API_KEY="change-me"
.\api\.venv\Scripts\uvicorn app.main:app \
  --app-dir api --host 127.0.0.1 --port 8000
\end{verbatim}

Use the same values as in \texttt{.env}. Keep this window open.

\textbf{PowerShell 2:}
\begin{verbatim}
cd C:\path\to\chess-club-manager
$env:CHESS_API_URL="http://127.0.0.1:8000"
$env:CHESS_API_KEY="change-me"
$env:PYTHONPATH="$PWD\frontend"
py -m chess_club_fe.main
\end{verbatim}

\subsection{Quick backend check}

Before opening the GUI, open the following address in a browser:

\begin{verbatim}
http://127.0.0.1:8000/health
\end{verbatim}

A working backend returns:

\begin{verbatim}
{"status":"ok"}
\end{verbatim}

\section{Using the Application}

\subsection{Players}

Open \textbf{Players} to view existing players. Click \textbf{Add Player} to create one. Enter a name, optional birth year and initial Elo. The default Elo is 1000. Creating a player also creates the first rating-history entry automatically.

\subsection{Create Tournament}

Open \textbf{Create Tournament}, enter a name, a start date in \texttt{YYYY-MM-DD} format and the number of rounds. The application creates the tournament and all required round records automatically.

After creation, the generated round IDs are shown, for example:

\begin{verbatim}
Rounds: R1=#1, R2=#2
\end{verbatim}

Keep these IDs because \textbf{Record Game} requires the database round ID.

\subsection{Register Player}

Open \textbf{Register Player}, choose a player and tournament, optionally enter a seed and click \textbf{Register}. Both players must be registered before their game can be recorded. A player cannot be registered twice for the same tournament.

\subsection{Record Game}

Open \textbf{Record Game} and enter the round ID, White player, Black player, ECO code and result. Valid results are \texttt{1-0}, \texttt{0-1}, and \texttt{1/2-1/2}.

A successful game immediately updates both Elo ratings and creates two new rating-history entries.

The supplied ECO reference data contains:

\begin{center}
\begin{tabular}{ll}
\toprule
\textbf{Code} & \textbf{Opening} \\
\midrule
A00 & Irregular Opening \\
B01 & Scandinavian Defense \\
C00 & French Defense \\
C50 & Italian Game \\
D02 & London System \\
E60 & King's Indian Defense \\
\bottomrule
\end{tabular}
\end{center}

\subsection{Standings and Rating History}

Open \textbf{Standings} and choose a tournament to see points, Buchholz, wins, draws, losses and games played.

Open \textbf{Rating History} to view Elo values over time. The dropdown can show all players or one selected player.

\section{Functional Verification}

The following short test checks the main functions from the user's perspective.

\begin{center}
\begin{tabular}{>{\raggedright\arraybackslash}p{3.4cm} >{\raggedright\arraybackslash}p{5.0cm} >{\raggedright\arraybackslash}p{5.0cm}}
\toprule
\textbf{Function} & \textbf{Test} & \textbf{Expected / Evaluation} \\
\midrule
Create player & Add Alice with Elo 1500 & Player appears and an initial history entry is created. \\
Create tournament & Create Club Open with one round & Tournament appears and one round ID is generated. \\
Register player & Register Alice and Bob & Both registrations succeed; duplicates are rejected. \\
Record game & Alice beats Bob using a valid round and ECO code & Game is saved and both Elo values change immediately. \\
Rating history & Open Alice in Rating History & Initial and updated Elo values are visible. \\
Standings & Open Club Open & Points, results and Buchholz are displayed. \\
\bottomrule
\end{tabular}
\end{center}

If these checks succeed, the complete core workflow is functioning correctly from the user's perspective.

\section{Troubleshooting and Scope}

\textbf{Cannot reach the API:} make sure FastAPI is running and check \texttt{http://127.0.0.1:8000/health}.

\textbf{Unauthorized / invalid API key:} make sure \texttt{CHESS\_API\_KEY} used by the frontend matches \texttt{API\_KEY} used by the backend.

\textbf{Docker command fails:} make sure Docker Desktop or Docker Engine is running.

\textbf{Player not registered:} register both players for the tournament before recording their game.

\textbf{Unknown round ID:} use one of the round IDs shown when the tournament was created.

\textbf{Unknown ECO code:} use one of the ECO codes listed in this manual.

The delivered version intentionally does not include automatic pairing generation, edit/delete operations through the GUI, automatic tournament closing, user accounts or detailed opening statistics.

\end{document}
